% Source PDF: source/普通高考/2014/2014天津理.pdf, page 1 (question 3).
% PDF embedded-bitmap attempt: pdfimages -list returned no image objects; vector flowchart.
% Node-edge list: 开始 -> S=1,i=1 -> T=2i+1 -> S=S×T -> i=i+1 -> i>=4?;
% 否 loops back to T=2i+1, while 是 -> 输出 S -> 结束.
% solid segments: all node borders and directed connectors; dashed segments: none.
% labels: node text is centered with zero paragraph indent; 否/是 are placed beside
% their corresponding decision branches. The flowchart is schematic, not to scale.
\begin{tikzpicture}[
    >=Stealth,
    line width=0.55pt,
    line cap=round,
    line join=round,
    flowtext/.style={anchor=center,align=center,
        execute at begin node={\setlength{\parindent}{0pt}}},
    every node/.style={flowtext,inner xsep=4pt,inner ysep=2pt},
    branchlabel/.style={flowtext,inner sep=0pt},
]
    \node[draw, rounded corners=5pt, minimum width=1.35cm, minimum height=0.64cm]
      (start) at (0,9.00) {开始};
    \node[draw, minimum width=2.85cm, minimum height=0.70cm]
      (init) at (0,7.55) {$S=1,i=1$};
    \node[draw, minimum width=2.20cm, minimum height=0.70cm]
      (term) at (0,6.05) {$T=2i+1$};
    \node[draw, minimum width=2.95cm, minimum height=0.70cm]
      (multiply) at (0,4.60) {$S=S\times T$};
    \node[draw, minimum width=2.25cm, minimum height=0.70cm]
      (increment) at (0,3.15) {$i=i+1$};
    \node[draw, diamond, aspect=2.65, inner sep=0pt,
      minimum width=3.75cm, minimum height=1.05cm]
      (check) at (0,1.60) {$i\geq 4?$};
    \node[draw, trapezium, trapezium left angle=72, trapezium right angle=108,
      minimum width=2.25cm, minimum height=0.76cm]
      (output) at (0,0.05) {输出 $S$};
    \node[draw, rounded corners=5pt, minimum width=1.35cm, minimum height=0.64cm]
      (finish) at (0,-1.35) {结束};

    \draw[->] (start.south) -- (init.north);
    \draw[->] (init.south) -- (term.north);
    \draw[->] (term.south) -- (multiply.north);
    \draw[->] (multiply.south) -- (increment.north);
    \draw[->] (increment.south) -- (check.north);

    % 否 branch: return to the main trunk above T=2i+1, with a separate
    % right-pointing merge arrow and the trunk's single downward arrow.
    \draw (check.west) -- node[branchlabel,above=2pt] {否}
      (-3.25,1.60) -- (-3.25,6.70);
    \draw[->] (-3.25,6.70) -- (0,6.70);
    % 是 branch proceeds independently to output and termination.
    \draw[->] (check.south) -- node[branchlabel,right=2pt] {是} (output.north);
    \draw[->] (output.south) -- (finish.north);
\end{tikzpicture}
